% Part: applied-modal-logics
% Chapter: epistemic-logic
% Section: language-epistemic-logic

\documentclass[../../../include/open-logic-section]{subfiles}

\begin{document}

\olfileid{aml}{el}{lan}

\olsection{જ્ઞાનસંબંધી તર્કની ભાષા}

\begin{defn}
$G$ એ કર્તા-સંકેતોનો ગણ હોય. બહુ-કર્તા જ્ઞાનસંબંધી તર્કની
મૂળભૂત ભાષામાં નીચેની બાબતો હોય છે:
\begin{enumerate}
  \tagitem{prvFalse}{!!{falsity} માટેનો વિધાનાત્મક અચલાંક~$\lfalse$.}{}
  \tagitem{prvTrue}{!!{truth} માટેનો વિધાનાત્મક અચલાંક~$\ltrue$.}{}
  \item !!{propositional variable}sનો !!{denumerable}s ગણ: $\Obj
    p_0$, $\Obj p_1$, $\Obj p_2$, \dots
  \item વિધાનાત્મક સંયોજકો: \startycommalist
  \iftag{prvNot}{\ycomma $\lnot$ (નિષેધ)}{}%
  \iftag{prvAnd}{\ycomma $\land$ (સંયોજન)}{}%
  \iftag{prvOr}{\ycomma $\lor$ (વિયોજન)}{}%
  \iftag{prvIf}{\ycomma $\lif$ (!!{conditional})}{}%
  \iftag{prvIff}{\ycomma $\liff$ (!!{biconditional})}{}.
  \item $a \in G$ હોય ત્યારે જ્ઞાન-સંક્રિયક $\Knows_a$.
\end{enumerate}
\end{defn}

જો આપણાં તંત્રમાં માત્ર એક જ કર્તાના જ્ઞાનનો વિચાર કરીએ, તો
ગણ~$G$ અને જુદા જુદા કર્તાઓનો ઉલ્લેખ છોડી શકીએ. એ કિસ્સામાં
માત્ર મૂળભૂત સંક્રિયક~$\Knows$ રહે છે.

\begin{defn}
જ્ઞાનસંબંધી ભાષાનાં \emph{!!^{formula}s} નીચે મુજબ આગમનાત્મક રીતે
  વ્યાખ્યાયિત થાય છે:
\begin{enumerate}
\tagitem{prvFalse}{$\lfalse$ પરમાણુક !!{formula} છે.}{}

\tagitem{prvTrue}{$\ltrue$ પરમાણુક !!{formula} છે.}{}

\item દરેક વિધાનાત્મક ચલ $\Obj p_i$ (પરમાણુક) !!{formula} છે.

\tagitem{prvNot}{જો $!A$ એ !!a{formula} હોય, તો $\lnot !A$ પણ
  !!a{formula} છે.}{}

\tagitem{prvAnd}{જો $!A$ અને $!B$ એ !!{formula}s હોય, તો $(!A \land
  !B)$ પણ !!a{formula} છે.}{}

\tagitem{prvOr}{જો $!A$ અને $!B$ એ !!{formula}s હોય, તો $(!A \lor !B)$
  પણ !!a{formula} છે.}{}

\tagitem{prvIf}{જો $!A$ અને $!B$ એ !!{formula}s હોય, તો $(!A \lif !B)$
  પણ !!a{formula} છે.}{}

\tagitem{prvIff}{જો $!A$ અને $!B$ એ !!{formula}s હોય, તો $(!A \liff !B)$
  પણ !!a{formula} છે.}{}

\item જો $!A$ એ !!a{formula} હોય અને $a \in G$ હોય, તો $\Knows_a !A$
  પણ !!a{formula} છે.

\tagitem{limitClause}{આ સિવાય બીજું કશું !!a{formula} નથી.}{}
\end{enumerate}
\end{defn}

જો !!a{formula}~$!A$માં $\Knows_a$ ન હોય, તો તેને
\emph{મોડલ-મુક્ત} કહીએ છીએ.

\begin{defn}
  સંક્રિયક~$\Knows$ વ્યક્તિગત જ્ઞાન દર્શાવવા માટે છે, જ્યારે
  $\EKnows$, જેને ઘણી વાર ``દરેક જણ જાણે છે'' એમ વાંચીએ છીએ,
  સમૂહ-જ્ઞાન દર્શાવે છે. $G' \subseteq G$ હોય ત્યારે
  $\EKnows_{G'} !A$ને \[\bigwedge_{b \in G'} \Knows_b !A.\]
  માટેના સંક્ષેપ તરીકે વ્યાખ્યાયિત કરીએ છીએ.
\end{defn}

જ્ઞાનનો એથી પણ વધુ પ્રબળ ખ્યાલ, એટલે કે કર્તાઓના ગણ~$G$માં
\emph{સામાન્ય જ્ઞાન}, પણ વ્યાખ્યાયિત કરી શકીએ છીએ. કોઈ માહિતી
કર્તાઓના સમૂહમાં સામાન્ય જ્ઞાન હોય, તો તેનો અર્થ એ છે કે એ સમૂહના
કર્તાઓના દરેક સંયોજન માટે સૌ જાણે છે કે બીજાં જાણે છે કે બીજાં
જાણે છે \dots એમ અનંત સુધી. આ સમૂહ-જ્ઞાન કરતાં ઘણું વધુ પ્રબળ
છે; એવા સંબંધાત્મક નિદર્શનો સરળતાથી રચી શકાય છે જેમાં
!!a{formula} સમૂહ-જ્ઞાન હોય પરંતુ સામાન્ય જ્ઞાન ન હોય.
``$G$માં~$!A$ સામાન્ય જ્ઞાન છે'' એ દર્શાવવા
$\CKnows_G !A$ વાપરીશું.

\end{document}
